% 396_LZ_Vorabfestlegung_De_ch.tex
% Johann Pascher, 28. Okt. 2026
% Vorab-Festlegung zum LZ-Ereignis bei 248 keV

\providecommand{\BM}{\textbf{[B]}}
\providecommand{\KM}{\textbf{[K]}}
\providecommand{\SM}{\textbf{[S]}}
\providecommand{\XM}{\textbf{[X]}}
\providecommand{\QM}{\textbf{[Q]}}

\chapter*{Vorab-Festlegung zum LZ-Ereignis bei 248 keV}
\addcontentsline{toc}{chapter}{Vorab-Festlegung zum LZ-Ereignis bei 248 keV}

\begin{abstract}
Der Xenon-Detektor LZ hat im September 2026 ein einzelnes Ereignis veröffentlicht, das nach allen
geprüften Kriterien wie der Stoß eines schweren Teilchens gegen einen Xenonkern aussieht, mit rund
248\,keV Rückstoßenergie, bei einem erwarteten Untergrund von etwa einem Hundertstel Ereignis. Die
globale Signifikanz beträgt 2{,}6\,$\sigma$; die Analyse ist nach Angabe der Kollaboration nicht
blind. In der FFGFT ist Dunkle Materie kein Stoff, sondern das Regime unterhalb der
Beschleunigungsskala $a_0$, in dem die Trägheit nichtlinear ist (Dok.~308). Ein direkter Nachweis
eines Dunkle-Materie-Teilchens ist deshalb ein Falsifikationskriterium dieser Lesart. Dieses
Dokument legt, bevor weitere Daten ausgewertet sind, fest, was die FFGFT für LZ erwartet und was sie
widerlegen würde: keine Häufung solcher Ereignisse und keine Jahresmodulation mit Maximum Anfang
Juni. Ein bestätigter Überschuss mit dieser Modulation würde die Lesart widerlegen. Der Ausweg, ein
solches Teilchen könne existieren, ohne die Rotationskurven zu tragen, wird ausdrücklich nicht in
Anspruch genommen.
\end{abstract}

% ============================================================
\section*{Statusmarkierungen}
% ============================================================
\begin{center}
\begin{tabular}{@{}lp{11cm}@{}}
\textbf{[B]} & Algebraisch bewiesen. \\[2pt]
\textbf{[K]} & Numerisch bestätigt (Prüfskript). \\[2pt]
\textbf{[S]} & Setzung, offen oder Einordnung ohne eigene Rechnung. \\[2pt]
\textbf{[X]} & Ausgeschlossen. \\[2pt]
\textbf{[Q]} & Aus der Literatur übernommen, mit Quelle. \\
\end{tabular}
\end{center}
Prüfskript: \texttt{2/python/Dok396\_Skripte/pruef\_396\_lz\_vorabfestlegung.py}.

% ============================================================
\section{Das Ereignis}
% ============================================================

LZ steht fast anderthalb Kilometer unter der Erde in der ehemaligen Homestake-Mine in South Dakota
und enthält sieben Tonnen flüssiges Xenon. Ein Stoß im Xenon erzeugt zwei Lichtblitze, einen im
Flüssigen und einen zweiten, wenn die freigesetzten Elektronen die Gasschicht erreichen. Aus dem
Verhältnis der beiden Blitze erkennt man, ob ein Elektron oder ein Atomkern gestoßen wurde; aus ihrer
Zeitdifferenz und der Lichtverteilung folgt der Ort des Stoßes.

Für die neue Auswertung hat die Kollaboration das Energiefenster für Kernrückstöße von $55$ auf
$270$\,keV erweitert und 220 Messtage mit 4{,}71 Tonnen Arbeitsmasse ausgewertet \QM~\cite{LZ2026}.
Darin liegt ein einziges Ereignis vom 16.~Juni 2023 mit $248\pm23\,(\text{stat})\pm23\,(\text{sys})$\,keV.
Es sitzt mitten im Arbeitsvolumen, mehr als 25\,cm von Wand und Boden entfernt, das äußere
Vetosystem hat nichts angezeigt, und es liegt genau im Band der Kernrückstöße, weit entfernt vom Band
der gewöhnlichen Elektronenstöße. Erwartet waren aus allen bekannten Untergrundquellen in diesem
Bereich $0{,}0106\pm0{,}0008$ Ereignisse \QM. Neutronen, ein Zerfall mit verschluckter Ladung,
Neutrinos und zufällige Überlagerungen haben die Autoren geprüft; keine dieser Quellen erklärt das
Ereignis überzeugend, vor allem fehlen die schwächeren Begleitstöße, die fast jede bekannte Quelle
hinterlassen müsste \QM.

Zwei Einschränkungen nennt die Kollaboration selbst. Erstens ist die Analyse nicht blind, weil die
zur Kontrolle beigemischten Kunstereignisse den Bereich hoher Energie nicht abdeckten. Zweitens
beträgt die Signifikanz nach Berücksichtigung der 616 geprüften Modellvarianten global
$2{,}6\,\sigma$, für einzelne Modelle lokal bis $3{,}4\,\sigma$ \QM. Ein Zufallsbild dieser Art
erscheint bei $2{,}6\,\sigma$ einmal in rund 215 Versuchen, für eine Entdeckung verlangt die
Teilchenphysik $5\,\sigma$, also einmal in rund $3{,}5$ Millionen \KM. Allein aus der Zählrate,
ohne die Lage des Ereignisses im Band, ergibt sich bei $0{,}0106$ erwarteten Ereignissen eine
Wahrscheinlichkeit von $1{,}05\,\%$ für mindestens eines \KM.

Unter den Deutungen mit Dunkler Materie passt am besten die inelastische Dunkle Materie
\QM~\cite{TuckerSmithWeiner2001}: Das Teilchen geht beim Stoß in einen etwas schwereren Zustand
über und bezahlt diese Energie aus seiner Bewegung. Langsame Teilchen können das nicht, deshalb
fehlen die schwachen Stöße, und es bleiben wenige kräftige. Für einen elastischen Stoß dieser Stärke
müsste das Teilchen mindestens etwa 75\,GeV schwer sein \KM.

% ============================================================
\section{Die Position der FFGFT}
% ============================================================

Die Befunde, die zur Dunklen Materie geführt haben, sind real: flache Rotationskurven,
Gravitationslinsen, der Bullet Cluster. Die FFGFT führt sie nicht auf einen zusätzlichen Stoff
zurück, sondern auf ein Kraftgesetz, das aus $\tilde T\cdot m=1$ zusammen mit der K3-Geometrie folgt,
\[
  a=\sqrt{a_N^2+a_N\,a_0},
\]
mit der Beschleunigungsskala $a_0$ (Dok.~308, 365). Dunkle Materie bezeichnet in dieser Lesart
keinen fehlenden Massenanteil, sondern das Regime unterhalb $a_0$, in dem die Trägheit nichtlinear
ist. Belegt ist das für die gasdominierte Zwerggalaxie DDO~154 \KM; für Spiralgalaxien bleiben
Abweichungen von rund 19\,\% offen, und für den Bullet Cluster ist es nicht belegt (Dok.~308, 384)
\SM.

Aus dieser Lesart folgt ohne weitere Annahme: Es gibt kein Teilchen der Dunklen Materie, also auch
keinen Halo aus solchen Teilchen, durch den sich die Erde bewegt. Das Buch \enquote{Dunkle Materie
-- eine Illusion} nennt deshalb den direkten Nachweis eines solchen Teilchens ausdrücklich als
Kriterium, an dem die Lesart scheitern würde. Das LZ-Ereignis ist genau eine Messung dieser Art.

% ============================================================
\section{Vorab-Festlegung}
% ============================================================

Die folgenden Aussagen gelten ab dem 28.~Oktober 2026, bevor die weiteren LZ-Daten ausgewertet sind.

\textbf{Erste Erwartung: keine Häufung.} In den noch nicht ausgewerteten Daten, einem Mehrfachen der
220 Messtage, gesammelt mindestens bis 2028, bleibt die Zahl der Kernrückstöße oberhalb von rund
$55$\,keV mit dem erwarteten Untergrund verträglich. Ein zweites oder drittes Ereignis bei einem
Untergrund von etwa $0{,}01$ je 220 Tage wäre bereits ein deutlicher Hinweis gegen diese Erwartung
\SM.

\textbf{Zweite Erwartung: keine Jahresmodulation.} Bewegt sich die Erde durch einen Halo aus
Dunkler Materie, addieren sich im Sommer Erd- und Sonnenbewegung. Die Geschwindigkeit gegen den Halo
schwankt dann um etwa $\pm6\,\%$ mit einem Maximum um den 2.~Juni \KM; bei inelastischer Dunkler
Materie, die nur die schnellsten Teilchen nachweist, wird daraus eine starke Schwankung der Rate.
Das bisherige Ereignis lag 14 Tage nach diesem Maximum, bei 97\,\% der vollen Phase \KM, was für
sich allein nichts besagt. Die FFGFT hat keinen Halo und erwartet daher keine solche Modulation \SM.

\textbf{Dritte Erwartung: keine unabhängige Bestätigung.} Andere Xenon-Detektoren, die denselben
Energiebereich untersuchen, finden keinen entsprechenden Überschuss \SM.

\textbf{Widerlegungskriterium.} Eine bestätigte Häufung solcher Ereignisse mit einem Jahresmaximum
Anfang Juni, oder eine unabhängige Bestätigung in einem zweiten Detektor mit mindestens
$5\,\sigma$, widerlegt die Lesart der Dunklen Materie als Regime ohne Stoff. Sie wäre dann mit \XM{}
zu kennzeichnen.

\textbf{Reichweite der Widerlegung.} Betroffen wäre nur diese eine Lesart: das Kraftgesetz mit
$a_0$ als alleinige Erklärung der Rotationskurven und Linseneffekte. Die übrigen Teile der FFGFT
hängen nicht an ihr. Die Herleitungen aus $\xi$ im Teilchensektor (Leptonmassen, $\alpha$, Koide-Form),
die Gravitation über $\tilde T\cdot m=1$ im Sonnensystem und die achromatische Rotverschiebung ohne
Expansion bleiben davon unberührt; die Rotverschiebungslesart wird an eigenen Kriterien geprüft
(Dok.~308). Ein Teilchen der Dunklen Materie wäre in der FFGFT dann eine zusätzliche Mode, deren
Herkunft aus der Struktur des Trägers erst zu zeigen wäre \SM.

\textbf{Kein nachträglicher Ausweg.} Ein einzelnes Ereignis legt nur das Produkt aus örtlicher
Dichte und Wirkungsquerschnitt fest. Man könnte deshalb einwenden, ein solches Teilchen könne
existieren, ohne die Rotationskurven zu tragen. Diesen Einwand nimmt die FFGFT nicht in Anspruch:
Er würde aus einer prüfbaren Lesart eine unprüfbare machen. Wer einen neuen Stoff zulässt, sobald er
gemessen wird, sagt über diesen Stoff vorher nichts.

% ============================================================
\section{Was das Ereignis nicht bedeutet}
% ============================================================

Ein einzelnes Ereignis bei $2{,}6\,\sigma$ widerlegt die FFGFT nicht, und es bestätigt sie auch
nicht. Es gibt Deutungen ohne Dunkle Materie, etwa ein seltenes Neutron, das die Abschirmung
umgangen und den Kern ohne Begleitstreuung getroffen hat, oder einen Vorgang im Xenonkern selbst
\QM. Über die Kernphysik des Xenons sagt die FFGFT nichts Eigenes; jede dieser Deutungen ist mit ihr
verträglich. Der Wert dieser Festlegung liegt darin, dass die FFGFT hier eine Erwartung hat, die von
der Standardkosmologie mit Teilchen-Halo abweicht und in den nächsten Jahren an Daten entschieden
wird, die heute noch niemand ausgewertet hat.

% ============================================================
\section{Übersicht}
% ============================================================
\begin{center}
\small
\begin{tabular}{@{}p{7.6cm}p{1.2cm}p{5.0cm}@{}}
\toprule
\textbf{Aussage} & \textbf{Status} & \textbf{Grund} \\
\midrule
Ereignis bei 248\,keV, Untergrund $0{,}0106$, $2{,}6\,\sigma$ global, nicht blind & \QM & LZ 2026 \\
$2{,}6\,\sigma$: etwa 1 zu 215; $5\,\sigma$: etwa 1 zu $3{,}5$ Mio. & \KM & Prüfskript \\
elastisch verlangt $m\gtrsim75$\,GeV & \KM & Kinematik \\
Dunkle Materie als Regime unterhalb $a_0$; DDO~154 & \KM & Dok.~308 \\
Spiralgalaxien, Bullet Cluster & \SM & Dok.~308, 384 \\
keine Häufung, keine Modulation, keine Bestätigung & \SM & Vorab-Festlegung \\
bestätigte Häufung mit Junimaximum widerlegt die Lesart & -- & Widerlegungskriterium \\
\bottomrule
\end{tabular}
\end{center}

\begin{thebibliography}{9}
\bibitem{LZ2026} LZ Collaboration, Suche nach Kernrückstößen bei höheren Energien im Rahmen der effektiven Feldtheorie der Dunklen Materie, Preprint, vorgestellt auf der TeVPA 2026 (1.~September 2026).
\bibitem{TuckerSmithWeiner2001} D.~Tucker-Smith, N.~Weiner, \textit{Inelastic dark matter}, Physical Review D \textbf{64} (2001) 043502.
\end{thebibliography}
